% !TEX TS-program = XeLaTeX
% !TEX encoding = UTF-8 Unicode

%%%%%%%%%%%%%%%%%%%%%%%%%%%%%%%%%%%%%%%%%%%%%%%%%%%%%%%%%%%%%%%%%%%%%%
%
% 大连理工大学硕士论文 XeLaTeX 模版 —— 格式配置文件 format.tex
% 版本：0.82
% 最后更新：2012.05.07
% 修改者：Yuri (E-mail: yuri_1985@163.com)
% 修订者：whufanwei(E-mail: dutfanwei@qq.com)
% 编译环境1：Ubuntu 12.04 + TeXLive 2011 + Emacs
% 编译环境2：Windows 7 + CTeX v2.9.2.164 + WinEdit
%
%%%%%%%%%%%%%%%%%%%%%%%%%%%%%%%%%%%%%%%%%%%%%%%%%%%%%%%%%%%%%%%%%%%%%%

%%%%%%%%%%%%%%%%%%%%%%%%%%%%%%%%%%%%%%%%%%%%%%%%%%%%%%%%%%%%%%%%%%%%%%
% 页面设置
%%%%%%%%%%%%%%%%%%%%%%%%%%%%%%%%%%%%%%%%%%%%%%%%%%%%%%%%%%%%%%%%%%%%%%
% A4 纸张
\setlength{\paperwidth}{21.0cm}
\setlength{\paperheight}{29.7cm}
% 设置正文尺寸大小
\setlength{\textwidth}{16.1cm}
\setlength{\textheight}{22.2cm}
% 设置正文区在正中间
\newlength \mymargin
\setlength{\mymargin}{(\paperwidth-\textwidth)/2}
\setlength{\oddsidemargin}{(\mymargin)-1in}
\setlength{\evensidemargin}{(\mymargin)-1in}
% 设置正文区偏移量，奇数页向右偏，偶数页向左偏
\newlength \myshift
\setlength{\myshift}{0.35cm}	% 双面打印的奇偶页偏移值，可根据需要修改，建议小于 0.5cm
\addtolength{\oddsidemargin}{\myshift}
\addtolength{\evensidemargin}{-\myshift}
% 页眉页脚相关距离设置
\setlength{\topmargin}{-0.05cm}
\setlength{\headheight}{0.50cm}
\setlength{\headsep}{0.90cm}
\setlength{\footskip}{1.47cm}
% 公式的精调
\allowdisplaybreaks[4]  % 可以让公式在排不下的时候分页排，这可避免页面有大段空白。

% 下面这组命令使浮动对象的缺省值稍微宽松一点，从而防止幅度
% 对象占据过多的文本页面，也可以防止在很大空白的浮动页上放置很小的图形。
\renewcommand{\topfraction}{0.9999999}
\renewcommand{\textfraction}{0.0000001}
\renewcommand{\floatpagefraction}{0.9999}

%%%%%%%%%%%%%%%%%%%%%%%%%%%%%%%%%%%%%%%%%%%%%%%%%%%%%%%%%%%%%%%%%%%%%%
% 字体字号定义
%%%%%%%%%%%%%%%%%%%%%%%%%%%%%%%%%%%%%%%%%%%%%%%%%%%%%%%%%%%%%%%%%%%%%%
% 字号
\newcommand{\yihao}{\fontsize{26pt}{39pt}\selectfont}	  % 一号，1.5  倍行距
\newcommand{\xiaoyi}{\fontsize{24pt}{30pt}\selectfont}  % 小一，1.25 倍行距
\newcommand{\erhao}{\fontsize{22pt}{27.5pt}\selectfont} % 二号，1.25 倍行距
\newcommand{\xiaoer}{\fontsize{18pt}{22.5pt}\selectfont}% 小二，1.25 倍行距
\newcommand{\sanhao}{\fontsize{16pt}{20pt}\selectfont}  % 三号，1.25 倍行距
\newcommand{\xiaosan}{\fontsize{15pt}{19pt}\selectfont} % 小三，1.25 倍行距
\newcommand{\sihao}{\fontsize{14pt}{17.5pt}\selectfont} % 四号，1.25倍行距
\newcommand{\xiaosi}{\fontsize{12pt}{15pt}\selectfont}  % 小四，1.25倍行距
\newcommand{\dawu}{\fontsize{10.5pt}{18pt}\selectfont}  % 五号，1.75倍行距
\newcommand{\zhongwu}{\fontsize{10.5pt}{16pt}\selectfont}% 五号，1.5 倍行距
\newcommand{\wuhao}{\fontsize{10.5pt}{10.5pt}\selectfont}% 五号，单倍行距
\newcommand{\xiaowu}{\fontsize{9pt}{9pt}\selectfont}	   % 小五，单倍行距
\newcommand{\xiaosandan}{\fontsize{15pt}{15pt}\selectfont}
\newcommand{\ershihao}{\fontsize{19.92pt}{19.92pt}\selectfont}
\newcommand{\xiaosiyidianwu}{\fontsize{12pt}{18pt}\selectfont}

\newcommand{\song}{\CJKfamily{song}}
\newcommand{\hei}{\CJKfamily{hei}}
\newcommand{\kai}{\CJKfamily{kai}}
\newcommand{\fs}{\CJKfamily{fs}}
\newcommand{\xkai}{\CJKfamily{xkai}}
\newcommand{\hwzsong}{\CJKfamily{hwzsong}}


% defaultfont 默认字体命令
\def\defaultfont{\renewcommand{\baselinestretch}{1.27}
  \fontsize{12pt}{15pt}\selectfont}

% 设置目录字体和行间距
\def\defaultmenufont{\renewcommand{\baselinestretch}{1.22}
  \fontsize{12pt}{15pt}\selectfont}

% 固定距离内容填入及下划线
\makeatletter
\newcommand\fixeddistanceleft[2][1cm]{{\hb@xt@ #1{#2\hss}}}
\newcommand\fixeddistancecenter[2][1cm]{{\hb@xt@ #1{\hss#2\hss}}}
\newcommand\fixeddistanceright[2][1cm]{{\hb@xt@ #1{\hss#2}}}
\newcommand\fixedunderlineleft[2][1cm]{\underline{\hb@xt@ #1{#2\hss}}}
\newcommand\fixedunderlinecenter[2][1cm]{\underline{\hb@xt@ #1{\hss#2\hss}}}
\newcommand\fixedunderlineright[2][1cm]{\underline{\hb@xt@ #1{\hss#2}}}
\makeatother

%%%%%%%%%%%%%%%%%%%%%%%%%%%%%%%%%%%%%%%%%%%%%%%%%%%%%%%%%%%%%%%%%%%%%%
% 标题环境相关
%%%%%%%%%%%%%%%%%%%%%%%%%%%%%%%%%%%%%%%%%%%%%%%%%%%%%%%%%%%%%%%%%%%%%%
% 定义、定理等环境
\theoremstyle{plain}
\theoremheaderfont{\hei\bf}
\theorembodyfont{\song\rmfamily}
\newtheorem{definition}{\hei 定义}[chapter]
\newtheorem{example}{\hei 例}[chapter]
\newtheorem{algorithm}{\hei 算法}[chapter]
\newtheorem{theorem}{\hei 定理}[chapter]
\newtheorem{axiom}{\hei 公理}[chapter]
\newtheorem{proposition}[theorem]{\hei 命题}
\newtheorem{property}{\hei 性质}
\newtheorem{lemma}[theorem]{\hei 引理}
\newtheorem{corollary}{\hei 推论}[chapter]
\newtheorem{remark}{\hei 注解}[chapter]
\newenvironment{proof}{\hei{证明}　}{\hfill $\square$ \vskip 4mm}

% 目录标题
\renewcommand\contentsname{\hfill 目　　录 \hfill}
\renewcommand\listfigurename{\hfill 插~图~目~录 \hfill}
\renewcommand\listtablename{\hfill 表~格~目~录 \hfill}
\renewcommand{\bibname}{\hfill 参~考~文~献 \hfill}

%%%%%%%%%%%%%%%%%%%%%%%%%%%%%%%%%%%%%%%%%%%%%%%%%%%%%%%%%%%%%%%%%%%%%%
% 段落章节相关
%%%%%%%%%%%%%%%%%%%%%%%%%%%%%%%%%%%%%%%%%%%%%%%%%%%%%%%%%%%%%%%%%%%%%%
\setcounter{secnumdepth}{4}
\setcounter{tocdepth}{4}
% 设置章、节、小节、小小节的间距
\titleformat{\chapter}[hang]{\normalfont\xiaosan\hei\sf}{\xiaosan\thechapter}{1em}{\xiaosan}
\titlespacing{\chapter}{0pt}{-3ex  plus .1ex minus .2ex}{3.3ex}
\titleformat{\section}[hang]{\sihao\hei\sf}{\sihao\thesection}{1em}{}{}
\titlespacing{\section}{0pt}{0.5em}{0.5em}
\titleformat{\subsection}[hang]{\xiaosi\hei\sf}{\xiaosi\thesubsection}{1em}{}{}
\titlespacing{\subsection}{0pt}{0.5em}{0.3em}
\titleformat{\subsubsection}[hang]{\hei\sf}{\thesubsubsection}{1em}{}{}
\titlespacing{\subsubsection}{0pt}{0.3em}{0pt}
% 缩小目录中各级标题之间的缩进
\dottedcontents{chapter}[0.32cm]{\vspace{0.2em}}{1.0em}{5pt}
\dottedcontents{section}[1.32cm]{}{1.8em}{5pt}
\dottedcontents{subsection}[2.32cm]{}{2.7em}{5pt}
\dottedcontents{subsubsection}[3.32cm]{}{3.4em}{5pt}

% 段落之间的竖直距离
\setlength{\parskip}{1.2pt}
% 段落缩进
\setlength{\parindent}{24pt}
% 定义行距
\renewcommand{\baselinestretch}{1.27}
% 参考文献条目间行间距
\setlength{\bibsep}{2pt}

%%%%%%%%%%%%%%%%%%%%%%%%%%%%%%%%%%%%%%%%%%%%%%%%%%%%%%%%%%%%%%%%%%%%%%
% 页眉页脚设置
%%%%%%%%%%%%%%%%%%%%%%%%%%%%%%%%%%%%%%%%%%%%%%%%%%%%%%%%%%%%%%%%%%%%%%

\newcommand{\makeheadrule}{%
  \makebox[0pt][l]{\rule[.7\baselineskip]{\headwidth}{0.5pt}}%
  \vskip-.8\baselineskip}

\makeatletter
\renewcommand{\headrule}{%
  {\if@fancyplain\let\headrulewidth\plainheadrulewidth\fi
    \makeheadrule}}

\pagestyle{fancyplain}

\fancyhf{}
\fancyhead[CO]{\song\wuhao{大连理工大学博士开题考核报告}}
\fancyhead[CE]{\song\wuhao\@ctitle}
\fancyfoot[C,C]{\xiaowu$-$~\thepage~$-$}

% Clear Header Style on the Last Empty Odd pages
\makeatletter
\def\cleardoublepagecover{\clearpage\if@twoside \ifodd\c@page\else%
  \hbox{}%
  \thispagestyle{empty}%              % Empty header styles
  \newpage%
  \if@twocolumn\hbox{}\newpage\fi\fi\fi}



%%%%%%%%%%%%%%%%%%%%%%%%%%%%%%%%%%%%%%%%%%%%%%%%%%%%%%%%%%%%%%%%%%%%%%
% 列表环境设置
%%%%%%%%%%%%%%%%%%%%%%%%%%%%%%%%%%%%%%%%%%%%%%%%%%%%%%%%%%%%%%%%%%%%%%
\setlist[enumerate]{(1),itemsep=-5pt,topsep=0mm,labelindent=\parindent,leftmargin=*}


%%%%%%%%%%%%%%%%%%%%%%%%%%%%%%%%%%%%%%%%%%%%%%%%%%%%%%%%%%%%%%%%%%%%%%
% 国际单位，以点连接。
%%%%%%%%%%%%%%%%%%%%%%%%%%%%%%%%%%%%%%%%%%%%%%%%%%%%%%%%%%%%%%%%%%%%%%
\sisetup{inter-unit-product = { }\cdot{ }}

%%%%%%%%%%%%%%%%%%%%%%%%%%%%%%%%%%%%%%%%%%%%%%%%%%%%%%%%%%%%%%%%%%%%%%
% 参考文献的处理
%%%%%%%%%%%%%%%%%%%%%%%%%%%%%%%%%%%%%%%%%%%%%%%%%%%%%%%%%%%%%%%%%%%%%%

% \addtolength{\bibsep}{-0.5em}              % 缩小参考文献间的垂直间距
\setlength{\bibhang}{2em}
\bibpunct{[}{]}{,}{s}{}{}



% \let\orig@Itemize =\itemize
% \let\orig@Enumerate =\enumerate
% \let\orig@Description =\description

% \def\Myspacing{\itemsep=1ex \topsep=-4ex \partopsep=-2ex \parskip=-1ex \parsep=2ex}
% \def\newitemsep{
% \renewenvironment{itemize}{\orig@Itemize\Myspacing}{\endlist}
% \renewenvironment{enumerate}{\orig@Enumerate\Myspacing}{\endlist}
% \renewenvironment{description}{\orig@Description\Myspacing}{\endlist}
% }
%   \def\olditemsep{
%   \renewenvironment{itemize}{\orig@Itemize}{\endlist}
%   \renewenvironment{enumerate}{\orig@Enumerate}{\endlist}
%   \renewenvironment{description}{\orig@Description}{\endlist}
% }
%   \renewcommand{\labelenumi}{(\arabic{enumi})}
%   \newitemsep

%%%%%%%%%%%%%%%%%%%%%%%%%%%%%%%%%%%%%%%%%%%%%%%%%%%%%%%%%%%%%%%%%%%%%%
%   其他设置
%%%%%%%%%%%%%%%%%%%%%%%%%%%%%%%%%%%%%%%%%%%%%%%%%%%%%%%%%%%%%%%%%%%%%%
%   增加 \ucite 命令使显示的引用为上标形式
%   \newcommand{\ucite}[1]{$^{\mbox{\scriptsize \cite{#1}}}$}

%%%%%%%%%%%%%%%%%%%%%%%%%%%%%%%%%%%%%%%%%%%%%%%%%%%%%%%%%%%%%%%%%%%%%%
%   图形表格
%%%%%%%%%%%%%%%%%%%%%%%%%%%%%%%%%%%%%%%%%%%%%%%%%%%%%%%%%%%%%%%%%%%%%%
\renewcommand{\figurename}{图}
\renewcommand{\tablename}{表}
% \captionstyle{\centering}
% \hangcaption
\captiondelim{\hspace{1em}}
\captionnamefont{\zhongwu}
\captiontitlefont{\zhongwu}
\setlength{\abovecaptionskip}{0pt}
\setlength{\belowcaptionskip}{0pt}


\newcommand{\tablepage}[2]{\begin{minipage}{#1}\vspace{0.5ex} #2 \vspace{0.5ex}\end{minipage}}
\newcommand{\returnpage}[2]{\begin{minipage}{#1}\vspace{0.5ex} #2 \vspace{-1.5ex}\end{minipage}}


%%%%%%%%%%%%%%%%%%%%%%%%%%%%%%%%%%%%%%%%%%%%%%%%%%%%%%%%%%%%%%%%%%%%%%
% 定义题头格言的格式
%%%%%%%%%%%%%%%%%%%%%%%%%%%%%%%%%%%%%%%%%%%%%%%%%%%%%%%%%%%%%%%%%%%%%%

\newsavebox{\AphorismAuthor}
\newenvironment{Aphorism}[1]
{\vspace{0.5cm}\begin{sloppypar} \slshape
    \sbox{\AphorismAuthor}{#1}
    \begin{quote}\small\itshape }
    {\\ \hspace*{\fill}------\hspace{0.2cm} \usebox{\AphorismAuthor}
    \end{quote}
  \end{sloppypar}\vspace{0.5cm}}

% 自定义一个空命令，用于注释掉文本中不需要的部分。
\newcommand{\comment}[1]{}

% This is the flag for longer version
\newcommand{\longer}[2]{#1}

\newcommand{\ds}{\displaystyle}

% define graph scale
\def\gs{1.0}

%%%%%%%%%%%%%%%%%%%%%%%%%%%%%%%%%%%%%%%%%%%%%%%%%%%%%%%%%%%%%%%%%%%%%%%%%%%%%%%%
% 封面摘要
%%%%%%%%%%%%%%%%%%%%%%%%%%%%%%%%%%%%%%%%%%%%%%%%%%%%%%%%%%%%%%%%%%%%%%%%%%%%%%%%
\def\cdegree#1{\def\@cdegree{#1}}\def\@cdegree{}
\def\ctitle#1{\def\@ctitle{#1}}\def\@ctitle{}
\def\caffil#1{\def\@caffil{#1}}\def\@caffil{}
\def\csubject#1{\def\@csubject{#1}}\def\@csubject{}
\def\cauthor#1{\def\@cauthor{#1}}\def\@cauthor{}
\def\cauthorno#1{\def\@cauthorno{#1}}\def\@cauthorno{}
\def\csupervisor#1{\def\@csupervisor{#1}}\def\@csupervisor{}
\def\cdate#1{\def\@cdate{#1}}\def\@cdate{}
\def\croom#1{\def\@croom{#1}}\def\@croom{}
\def\cmytype#1{\def\@cmytype{#1}}\def\@cmytype{}
\def\cmydate#1{\def\@cmydate{#1}}\def\@cmydate{}
\def\studid#1{\def\@studid{#1}}\def\@studid{}

\long\def\cabstract#1{\long\def\@cabstract{#1}}\long\def\@cabstract{}
\def\ckeywords#1{\def\@ckeywords{#1}}\def\@ckeywords{}
\def\etitle#1{\def\@etitle{#1}}\def\@etitle{}
\long\def\eabstract#1{\long\def\@eabstract{#1}}\long\def\@eabstract{}
\def\ekeywords#1{\def\@ekeywords{#1}}\def\@ekeywords{}

% 封面
\def\makecover{
  \begin{titlepage}
    \newpage
    \thispagestyle{empty}
    \begin{center}
    
    \begin{figure}[htbp]
    \centering
    \includegraphics[scale=1]{DUT}
    \end{figure}

    \parbox[t][4.8cm][t]{\textwidth}{
    \begin{center}\xiaoyi\hwzsong \textbf{博~士~生~开~题~考~核~报~告}\end{center} }

    \parbox[t][2.3cm][t]{\textwidth}{
    \begin{center}\sanhao\song
        论文题目：\makebox[10cm]\@ctitle
    \end{center} }

     \parbox[t][10.8cm][b]{\textwidth}
     {\xiaosandan
    \begin{center} \renewcommand{\arraystretch}{1.8} \song
    \begin{tabular}{lc@{：}cl}
    ~& { 姓\hfill 名} & \underline{\makebox[8.08cm][c]\@cauthor} &\\
    ~& { 学\hfill 号} & \underline{\makebox[8.08cm][c]\@studid} &\\
    ~& { 学\hfill 科\hfill 领\hfill 域} & \underline{\makebox[8.08cm][c]\@csubject} &\\
    ~& { 培\hfill 养\hfill 类\hfill 型} & \underline{\xiaosi\makebox[8.08cm][l]\@cmytype} &\\

    ~& { 指\hfill 导\hfill 教\hfill 师} & \underline{\makebox[8.08cm][c]\@csupervisor}&\\
    ~& { 入\hfill 学\hfill 日\hfill 期} & \underline{\makebox[8.08cm][c] \@cmydate}&\\
    ~& { 报\hfill 告\hfill 日\hfill 期} & \underline{\makebox[8.08cm][c]\@cdate}&\\
    ~& { 报\hfill 告\hfill 地\hfill 点} & \underline{\makebox[8.08cm][c]\@croom}&
    \end{tabular} \renewcommand{\arraystretch}{1}
    \end{center}}

    \vspace{0.7cm}
    \parbox[t][1cm][t]{\textwidth}{
    \begin{center}\sanhao\song{研究生院制表}\end{center} }
    \end{center}
    % \cleardoublepagecover
  \end{titlepage}
}

\def\Declare 
{
  \renewcommand{\baselinestretch}{1.61}
  % \newpage
  \thispagestyle{empty}
  \begin{center}
    % \parbox[t][1.52cm][c]{\textwidth}
    % {\erhao\kai\centerline{说~~~~明}}
    \erhao\kai~说~~~~明~
    % \parbox[t][12.68cm][c]{\textwidth}
    \end{center}
    % {
      \xiaosiyidianwu\song\noindent

      博士生的开题考核是博士生课程学习结束后开展课题研究工作的基本要求，是保证学位论文质量，工作进度和博士研究生培养质量的首要环节。%
      
           一、考核内容：首先，考查博士生对本学科专业的基础理论与专业知识的掌握程度、课程学习情况、是否具备应用所学知识进行创造性科研工作能力；%
      其次，考查博士学位论文工作准备情况，包括论文选题、文献阅读、工作难度、研究思路、预期成果与创新性、研究基础、写作能力和答辩表达能力等；%
      此外，还要考查博士生参与学术活动情况、政治思想素质、及学习和工作态度等。%

          二、报告撰写：开题考核报告正文字数不少于15000字；参考文献数量不少于50篇，其中，外文资料不少于二分之一，近5年文献不少于三分之一；
      正文及参考文献等撰写要求参见《大连理工大学博士学位论文格式规范》。

          三、报告保存：博士生根据专家的意见修改完善开题考核报告，经导师和点长签字生成PDF文档，上传到研究生信息管理系统，导师审核，由学部（学院）进行网上确认。
      
    % }
    
  
  \renewcommand{\baselinestretch}{1.27}
  \cleardoublepagecover
}


\def\makeabstract{
  \defaultfont
  \chapter*{\hfill 摘　　要 \hfill}
  \addcontentsline{toc}{chapter}{摘　　要}
  \setcounter{page}{1}
  \@cabstract
  \vspace{0.53cm}

  \noindent {\hei{关键词：{\fs\@ckeywords}}}

  \defaultfont
  \cleardoublepage
  \chapter*{}
  \addcontentsline{toc}{chapter}{Abstract}
  \vspace{-1.40cm}
  \begin{center}
    {\xiaosan\textrm{\@etitle}}
  \end{center}
  \vspace{-0.35cm}
  \begin{center}
    {
      \xiaosan{Abstract}\\
    }
  \end{center}
  \vspace{0.12cm}

  \@eabstract

  \vspace{0.55cm}

  \noindent {\textbf{Key Words:}}~~{\textsf{\@ekeywords}}
  \cleardoublepage
}

\makeatletter
\def\hlinewd#1{%
  \noalign{\ifnum0=`}\fi\hrule \@height #1 \futurelet
  \reserved@a\@xhline}
\makeatother

% 定义索引生成
\def\generateindex
{
  \addcontentsline{toc}{chapter}{\indexname}
  \printindex
  \cleardoublepage
}

\raggedbottom
